\documentclass[uplatex,dvipdfmx,a4paper,11pt]{jsarticle}
\input{../../../templates/note.tex}
\noteprint{統計力学2　第8章「縮退した気体と粒子の統計性」}{2026-11-25}
%% note-moto	problems\統計力学2\shou08\統計力学2-8.tex
\begin{document}
\begin{kirokuichiran}
  \ichirangyou{\hyperlink{bt-m-001-levels}{1}}{L1・計算}{○◎△}{B}{2026-10-25 (1)記述 (3)計算}
  \ichirangyou{\hyperlink{bt-m-002-lambda}{2}}{L1・計算}{○◎◎}{A}{2026-10-25 (1)単位}
  \ichirangyou{3}{L1・計算}{}{未着手}{}
  \ichirangyou{4}{L1・計算}{}{未着手}{}
  \ichirangyou{5}{L1・計算}{}{未着手}{}
  \ichirangyou{6}{L1・計算}{}{未着手}{}
  \ichirangyou{7}{L1・計算}{}{未着手}{}
  \ichirangyou{\hyperlink{bt-m-003-z1}{8}}{L2・導出}{◎◎◎◎}{S}{2026-10-25}
  \ichirangyou{\hyperlink{bt-m-004-entropy}{9}}{L2・計算}{◎△◎}{B}{2026-10-25 (2)計算}
  \ichirangyou{\hyperlink{bt-m-005-senyuu}{10}}{L2・計算}{○○◎○}{B}{2026-10-26 (1)記述 (2)単位 (4)記述}
  \ichirangyou{11}{L2・導出}{}{未着手}{}
  \ichirangyou{12}{L2・計算}{}{未着手}{}
  \ichirangyou{13}{L2・導出}{}{未着手}{}
  \ichirangyou{14}{L2・計算}{}{未着手}{}
  \ichirangyou{15}{L2・計算}{}{未着手}{}
  \ichirangyou{16}{L3・導出}{}{未着手}{}
  \ichirangyou{17}{L3・説明}{}{未着手}{}
  \ichirangyou{18}{L3・導出}{}{未着手}{}
  \ichirangyou{19}{L3・計算}{}{未着手}{}
  \ichirangyou{20}{L4・導出}{}{未着手}{}
\end{kirokuichiran}

\mondaihen
\level{1}
% toubun: kiroku\toubun\統計力学2\001_2026-10-25.tex
\mondaibangou{1}
\begin{mondai}{001-levels}{計算}{箱の中の1粒子のエネルギー準位と縮退度，$\beta E_0$ と $\lambda/L$ の関係，和を積分に置き換える根拠}
\begin{teisuu}
  換算プランク定数 & $\hbar$ & \qty{1.05e-34}{\joule\second}\\
  ボルツマン定数 & $\kB$ & \qty{1.38e-23}{\joule\per\kelvin}\\
  ヘリウム4原子の質量 & $m_{\mathrm{He}}$ &\qty{6.65e-27}{\kilogram}\\
  電子の質量 & $\melec$ & \qty{9.11e-31}{\kilogram}\\
\end{teisuu}
一辺 $L$ の立方体の容器に閉じ込められた質量 $m$ の1粒子のエネルギー固有値は，量子数 $n_x,n_y,n_z$（それぞれ $1,2,3,\dots$）を用いて
  \[ E_{n_xn_yn_z}=E_0\bigl(n_x^{2}+n_y^{2}+n_z^{2}\bigr),\qquad E_0=\frac{\pi^{2}\hbar^{2}}{2mL^{2}} \]
  で与えられる。温度を $T$，$\beta=1/(\kB T)$ とする。
  \begin{shoumon}
    \item エネルギーの値を低いほうから異なるものを5つ挙げ，それぞれを $E_0$ を単位として表せ。また，各値をとる1粒子状態の数（縮退度）を答えよ。
    \item 熱的ドブロイ波長 $\lambda=\sqrt{2\pi\hbar^{2}\beta/m}$ を用いて，$\beta E_0$ を $\lambda/L$ で表せ。
    \item $\beta E_0\ll1$ のとき，最低の準位に対する2番目の準位のボルツマン因子（1つの状態あたりの $\ee^{-\beta E}$）の比を $\beta E_0$ の1次まで求めよ。また，1粒子の分配関数の和を積分に置き換えてよい理由を述べよ。
  \end{shoumon}
\notekiroku{2026-10-25 B ○◎△}
\begin{kotaeutsushi}{2026-10-25}
\toi{1}{○}{記述}
$\{n_x,n_y,n_z\}$

i) $=\{1,1,1\}$

$E=3E_0$　縮退：1

ii) $=\{2,1,1\}\{1,2,1\}\{1,1,2\}$

$E=6E_0$　縮退：3

iii) $=\{2,2,1\}\{2,1,2\}\ayamari{\{1,1,2\}}{\{1,2,2\}}$

$E=9E_0$　縮退：3
\naze{$\{1,1,2\}$ は2乗の和が $6$ で，ii) の $6E_0$ の状態である。$9E_0$ の3つ目は $(2,2,1)$ の並べ替えの $\{1,2,2\}$ である。縮退度3は正しく数えている。}

iv) $=\{3,1,1\}\{1,3,1\}\{1,1,3\}$

$E=11E_0$　縮退：3

v) $=\{2,2,2\}$

$E=12E_0$　縮退：1
\toi{2}{◎}{}
\[ \beta E_0=\frac{\pi}{4}\cdot\frac{2\pi\hbar^{2}\beta}{m}\cdot\frac{1}{L^{2}}=\frac{\pi}{4}\Bigl(\frac{\lambda}{L}\Bigr)^{2} \]
\toi{3}{△}{計算}
\[ \frac{\ee^{-6\beta E_0}}{\ee^{-3\beta E_0}}=\ee^{-3\beta E_0} \]
\begin{align*}
\ee^{x}&=\frac{\ee^{0}}{0!}x^{0}+\ayamari{\frac{\ee^{1}}{1!}}{\frac{\ee^{0}}{1!}}x^{1}+O(x^{2})\\
&=1+\ee x+O(x^{2})
\end{align*}
\[ \ee^{-3\beta E_0}=1-3\ee\beta E_0 \]
\naze{テイラー展開の係数は $x=0$ での微分係数 $f^{(n)}(0)/n!$ で，$\ee^{x}$ の1階微分の $x=0$ での値は $\ee^{0}=1$ である。よって $\ee^{x}=1+x+O(x^{2})$ で，比は $1-3\beta E_0$ になる。以降はこの誤りを引き継ぐ。}

隣り合う準位のボルツマン因子の比がほぼ1であり，ゆっくりと変化するから。
\end{kotaeutsushi}
\begin{kizuki}
\item テイラー展開 $f(x)=\sum f^{(n)}(0)x^{n}/n!$ の係数は「$x=0$ で」の微分係数である。$\ee^{x}$ は何回微分しても $\ee^{x}$ だが，$x=0$ を入れるので係数はすべて $1$ になり，$\ee^{x}=1+x+x^{2}/2+\cdots$ である。$\ee^{1}=\ee$ は $x=1$ での値で，展開の係数には出てこない。
\item 確かめ方：近似式で $\beta E_0\to0$ とすると比は $1$ に近づくはずで，$1-3\ee\beta E_0$ もそうなるので1次の係数は別に確かめる。$\ee^{x}\simeq1+x$ は $x=0.01$ で $\ee^{0.01}=1.01005$ と合う（$1+\ee x=1.027$ は合わない）と数値で1回確かめる。
\item 状態の列は並べ替えで作る。$(2,2,1)$ の並べ替えは「$1$ をどこに置くか」で $\{1,2,2\}\{2,1,2\}\{2,2,1\}$ の3つと書き出し，各組の2乗の和が $9$ になるかを見て確かめる。
\end{kizuki}
\begin{tadashii}
\houshin (1) は $n_x^{2}+n_y^{2}+n_z^{2}$ の値を小さい順に列挙し，並べ替えの数を数える。(2) は $\lambda$ の定義を $\beta E_0$ に代入する。(3) はボルツマン因子の比を展開し，和の項が $n$ についてゆっくり変わることを示す。
\keisan \shoumonkai{1} $n_x,n_y,n_z\ge1$ なので，2乗の和 $s=n_x^{2}+n_y^{2}+n_z^{2}$ がとりうる値は小さい順に次のとおりである。
\begin{itemize}
  \item $s=3$：$(1,1,1)$ の1通り。
  \item $s=6$：$(2,1,1)$ の並べ替えで3通り。
  \item $s=9$：$(2,2,1)$ の並べ替えで3通り。
  \item $s=11$：$(3,1,1)$ の並べ替えで3通り。
  \item $s=12$：$(2,2,2)$ の1通り。
\end{itemize}
$s=4,5,7,8,10$ は，各項が $1,4,9,\dots$ のどれかである3項の和では作れない（例えば $10=9+1+0$ や $8=4+4+0$ は $0$ を含む）。次の値は $s=14$（$(3,2,1)$ の6通り）である。
\shoumonkai{2} $\lambda^{2}=2\pi\hbar^{2}\beta/m$ より $\hbar^{2}\beta/m=\lambda^{2}/(2\pi)$ なので
\[ \beta E_0=\frac{\pi^{2}}{2L^{2}}\cdot\frac{\hbar^{2}\beta}{m}=\frac{\pi^{2}}{2L^{2}}\cdot\frac{\lambda^{2}}{2\pi}=\frac{\pi}{4}\Bigl(\frac{\lambda}{L}\Bigr)^{2} \]
\shoumonkai{3} 最低の準位は $3E_0$，2番目は $6E_0$ なので，ボルツマン因子の比は
\[ \frac{\ee^{-6\beta E_0}}{\ee^{-3\beta E_0}}=\ee^{-3\beta E_0}\simeq1-3\beta E_0 \]
である。$\beta E_0\ll1$ ではこの比はほぼ $1$ で，隣り合う準位の寄与はほとんど変わらない。1つの量子数についての和の項 $f(n)=\ee^{-\beta E_0n^{2}}$ が大きく減るのは $n\sim1/\sqrt{\beta E_0}\gg1$ のときなので，$f(n)$ は $n$ を $1$ ずつ変えてもゆっくりとしか変わらない。したがって $n$ を連続変数とみなし，和を積分に置き換えてよい。
\kotae{\shoumonkai{1} $3E_0$（縮退度1），$6E_0$（3），$9E_0$（3），$11E_0$（3），$12E_0$（1） \shoumonkai{2} $\beta E_0=(\pi/4)(\lambda/L)^{2}$ \shoumonkai{3} 比は $1-3\beta E_0$。$f(n)$ が $n\sim1/\sqrt{\beta E_0}\gg1$ の幅でしか変わらないから}
\kenzan (1) 縮退度は並べ替えの数 $3!/(\text{同じ数字の個数の階乗の積})$ と一致する（$(2,1,1)$ は $3!/2!=3$，$(3,2,1)$ は $3!=6$）。(2) $\lambda/L$ は無次元で，$\beta E_0$ も無次元である。(3) $\beta E_0\to0$（高温または大きな容器）で比は $1$ に近づき，連続近似がよくなる。
\tsumazuki 量子数は $1$ から始まる。$0$ を許すと最低エネルギーを $0$ と誤る。
\end{tadashii}
\end{mondai}

% toubun: kiroku\toubun\統計力学2\002_2026-10-25.tex
\mondaibangou{2}
\begin{mondai}{002-lambda}{計算}{熱的ドブロイ波長，$\lambda^{3}\rho$ と $\beta E_0$ の数値評価による古典的な扱いの判定}
\begin{teisuu}
  換算プランク定数 & $\hbar$ & \qty{1.05e-34}{\joule\second}\\
  ボルツマン定数 & $\kB$ & \qty{1.38e-23}{\joule\per\kelvin}\\
  ヘリウム4原子の質量 & $m_{\mathrm{He}}$ &\qty{6.65e-27}{\kilogram}\\
  電子の質量 & $\melec$ & \qty{9.11e-31}{\kilogram}\\
\end{teisuu}
ヘリウム4原子の気体を，相互作用の無い同種粒子の気体とみなす。温度は $T=\qty{300}{\kelvin}$，数密度は $\rho=\qty{2.45e25}{\per\metre\cubed}$（ほぼ1気圧に相当）とする。定数は定数表の値を用い，答えは有効数字3桁で求めよ。
  \begin{shoumon}
    \item 熱的ドブロイ波長 $\lambda=\sqrt{2\pi\hbar^{2}\beta/m}$ を求めよ。
    \item $\lambda^{3}\rho$ を求め，この気体を古典的に（$Z_N=Z_1^{N}/N!$ で）扱ってよいかを判定せよ。
    \item 気体が一辺 $L=\qty{1.00}{\centi\metre}$ の立方体の容器に入っているとして $\beta E_0$ を求め，和を積分に置き換える仮定 $\beta E_0\ll1$ が成り立つかを判定せよ。
  \end{shoumon}
\notekiroku{2026-10-25 A ○◎◎}
\begin{kotaeutsushi}{2026-10-25}
\toi{1}{○}{単位}
\[ \lambda=\num{5.02e-11}\,\nuke{\unit{\metre}} \]
\naze{$\lambda$ は長さなので単位 $\unit{\metre}$ が要る。$\lambda=\qty{5.02e-11}{\metre}$ である。}
\toi{2}{◎}{}
\[ \lambda^{3}\rho=\num{3.09e-6}\ll1 \]
古典的に扱ってよい。
\toi{3}{◎}{}
\begin{align*}
\beta E_0&=\frac{\pi}{4}\Bigl(\frac{\lambda}{L}\Bigr)^{2}\\
&=\num{1.98e-17}\ll1
\end{align*}
\end{kotaeutsushi}
\begin{kizuki}
\item 数値の答えには単位を付ける。$\lambda$ は長さ（$\unit{\metre}$），$\lambda^{3}\rho$ と $\beta E_0$ は無次元なので単位なし，と答えを書く前に次元を1回確かめる。
\end{kizuki}
\begin{tadashii}
\houshin $\lambda$ の定義に数値を代入し，$\lambda^{3}\rho$ と $\beta E_0=(\pi/4)(\lambda/L)^{2}$ を計算して，それぞれ $1$ と比べる。途中の値は丸めずに計算する。
\keisan \shoumonkai{1} $\beta=1/(\kB T)$ を代入すると
\[ \lambda^{2}=\frac{2\pi\hbar^{2}}{m_{\mathrm{He}}\kB T}=\frac{2\pi\times(\num{1.05e-34})^{2}}{\num{6.65e-27}\times\num{1.38e-23}\times\num{300}}\,\unit{\metre\squared}=\frac{\num{6.927e-68}}{\num{2.753e-47}}\,\unit{\metre\squared}=\qty{2.516e-21}{\metre\squared} \]
なので $\lambda=\qty{5.016e-11}{\metre}$ である。
\shoumonkai{2} $\lambda^{3}=\qty{1.262e-31}{\metre\cubed}$ より
\[ \lambda^{3}\rho=\num{1.262e-31}\times\num{2.45e25}=\num{3.09e-6} \]
これは $1$ よりはるかに小さいので，1つの1粒子状態の平均の占有数は $1$ よりはるかに小さく，古典的に扱ってよい。
\shoumonkai{3} $L=\qty{1.00e-2}{\metre}$ より $(\lambda/L)^{2}=\num{2.516e-17}$ なので
\[ \beta E_0=\frac{\pi}{4}\Bigl(\frac{\lambda}{L}\Bigr)^{2}=\num{0.7854}\times\num{2.516e-17}=\num{1.98e-17} \]
これは $1$ よりはるかに小さいので，和を積分に置き換える仮定は十分に成り立つ。
\kotae{\shoumonkai{1} $\lambda=\qty{5.02e-11}{\metre}$ \shoumonkai{2} $\lambda^{3}\rho=\num{3.09e-6}\ll1$，古典的に扱ってよい \shoumonkai{3} $\beta E_0=\num{1.98e-17}\ll1$，仮定は成り立つ}
\kenzan 次元：$\hbar^{2}/(m\kB T)$ の単位は $\unit{\joule\squared\second\squared}/(\unit{\kilogram\joule})=\unit{\joule\second\squared\per\kilogram}=\unit{\metre\squared}$ で，$\lambda$ は長さである。オーダー：$\lambda\simeq\qty{0.05}{\nano\metre}$ は原子の大きさより小さく，平均の粒子間距離 $\rho^{-1/3}\simeq\qty{3.4}{\nano\metre}$ よりはるかに短いので，$\lambda^{3}\rho\ll1$ と整合する。
\tsumazuki $\lambda$ の定義の $2\pi$ を落としたり，$\hbar$ の代わりに $h$ を使ったりしないこと。
\end{tadashii}
\end{mondai}

\level{2}
% toubun: kiroku\toubun\統計力学2\003_2026-10-25.tex
\mondaibangou{8}
\begin{mondai}{003-z1}{導出}{1粒子の分配関数の計算（和の分解・和から積分への置き換え・ガウス積分），$Z_N$ からの内部エネルギーと自由エネルギー}
\begin{teisuu}
  換算プランク定数 & $\hbar$ & \qty{1.05e-34}{\joule\second}\\
  ボルツマン定数 & $\kB$ & \qty{1.38e-23}{\joule\per\kelvin}\\
  ヘリウム4原子の質量 & $m_{\mathrm{He}}$ &\qty{6.65e-27}{\kilogram}\\
  電子の質量 & $\melec$ & \qty{9.11e-31}{\kilogram}\\
\end{teisuu}
一辺 $L$（体積 $V=L^{3}$）の立方体の容器の中に，質量 $m$ の $N$ 個の同種粒子（相互作用なし）がある。1粒子のエネルギーは $E_\alpha=E_0(n_x^{2}+n_y^{2}+n_z^{2})$（$n_x,n_y,n_z=1,2,\dots$，$E_0=\pi^{2}\hbar^{2}/(2mL^{2})$）で，分配関数は $Z_N=Z_1^{N}/N!$ とする。温度を $T$，$\beta=1/(\kB T)$，$\lambda=\sqrt{2\pi\hbar^{2}\beta/m}$，$\rho=N/V$ とする。
  \begin{shoumon}
    \item $Z_1$ を，1つの量子数 $n$ についての和の3乗で表せ。
    \item $\beta E_0\ll1$ として和を積分に置き換え（積分の下限は $0$ としてよい），ガウス積分
    \[ \int_{0}^{\infty}\ee^{-x^{2}}\dd x=\frac{\sqrt{\pi}}{2} \]
    を用いて $Z_1=V/\lambda^{3}$ を示せ。
    \item スターリングの公式 $\log N!\simeq N\log N-N$ を用いて，$\log Z_N$ を $N$ と $\lambda^{3}\rho$ で表せ。
    \item $E=-\pdv{}{\beta}\log Z_N$ を用いて $E=(3/2)N\kB T$ を示せ。また，ヘルムホルツの自由エネルギー $F=-\kB T\log Z_N$ を $N$，$\kB T$，$\lambda^{3}\rho$ で表せ。
  \end{shoumon}
\notekiroku{2026-10-25 S ◎◎◎◎}
\begin{kotaeutsushi}{2026-10-25}
\toi{1}{◎}{}
\begin{align*}
Z_1&=\sum_{n_x=1}^{\infty}\sum_{n_y=1}^{\infty}\sum_{n_z=1}^{\infty}\ee^{-\beta E_0(n_x^{2}+n_y^{2}+n_z^{2})}\\
&=\biggl(\sum_{n=1}^{\infty}\ee^{-\beta E_0n^{2}}\biggr)^{3}
\end{align*}
\toi{2}{◎}{}
\begin{align*}
Z_1&=\biggl(\int_{0}^{\infty}\ee^{-\beta E_0x^{2}}\dd x\biggr)^{3}\\
&=\biggl(\frac{\sqrt{\pi}}{2\sqrt{\beta E_0}}\biggr)^{3}=\Bigl(\frac{\pi}{4\beta E_0}\Bigr)^{3/2}\\
&=\Bigl(\frac{2\pi mL^{2}}{4\pi^{2}\hbar^{2}\beta}\Bigr)^{3/2}=\Bigl(\frac{mL^{2}}{2\pi\hbar^{2}\beta}\Bigr)^{3/2}\\
&=\biggl(\sqrt{\frac{m}{2\pi\hbar^{2}\beta}}\biggr)^{3}L^{3}=\frac{V}{\lambda^{3}}
\end{align*}
\toi{3}{◎}{}
\begin{align*}
\log Z_N&=N\log Z_1-\log N!\\
&=N\log\frac{V}{\lambda^{3}}-N\log N+N\\
&=N\log\frac{V}{\lambda^{3}N}+N\\
&=-N\log\lambda^{3}\rho+N\\
&=N(1-\log\lambda^{3}\rho)
\end{align*}
\toi{4}{◎}{}
\begin{align*}
E&=-\pdv{}{\beta}\bigl(N(1-\log\lambda^{3}\rho)\bigr)\\
&=-\pdv{}{\beta}\Bigl(-\frac{3}{2}N\log\beta\Bigr)\\
&=\frac{3}{2}\cdot\frac{N}{\beta}\\
&=\frac{3}{2}N\kB T
\end{align*}
\begin{align*}
F&=-\kB TN(1-\log\lambda^{3}\rho)\\
&=\kB TN(\log\lambda^{3}\rho-1)
\end{align*}
\end{kotaeutsushi}
\end{mondai}

% toubun: kiroku\toubun\統計力学2\004_2026-10-25.tex
\mondaibangou{9}
\begin{mondai}{004-entropy}{計算}{エントロピーを $\lambda^{3}\rho$ で表すこと，$S<0$ となる条件，古典的な扱いが破れる温度の評価}
\begin{teisuu}
  換算プランク定数 & $\hbar$ & \qty{1.05e-34}{\joule\second}\\
  ボルツマン定数 & $\kB$ & \qty{1.38e-23}{\joule\per\kelvin}\\
  ヘリウム4原子の質量 & $m_{\mathrm{He}}$ &\qty{6.65e-27}{\kilogram}\\
  電子の質量 & $\melec$ & \qty{9.11e-31}{\kilogram}\\
\end{teisuu}
$N$ 個の同種粒子（質量 $m$，相互作用なし）からなる体積 $V$ の理想気体を，$Z_N=Z_1^{N}/N!$，$\beta E_0\ll1$，スターリングの公式で扱う。このとき
  \[ \log Z_N=N\bigl[1-\log(\lambda^{3}\rho)\bigr],\qquad E=\frac{3}{2}N\kB T \]
  である（$\lambda=\sqrt{2\pi\hbar^{2}\beta/m}$，$\rho=N/V$，$\beta=1/(\kB T)$）。これを用いてよい。
  \begin{shoumon}
    \item $S=(E-F)/T$ と $F=-\kB T\log Z_N$ を用いて，$S/(N\kB)$ を $\lambda^{3}\rho$ で表せ。
    \item $S<0$ となる条件を $\lambda^{3}\rho$ で表せ。さらに $\beta$ を消去して，その条件を $E$，$N$，$\rho$，$m$，$\hbar$ で表せ。
    \item 液体ヘリウム程度の数密度 $\rho=\qty{2.18e28}{\per\metre\cubed}$ のヘリウム4原子の集団を，この理想気体の式で扱うとする。$\lambda^{3}\rho=1$ となる温度 $T_1$ と，(1) で求めた $S$ が $0$ になる温度 $T_0$ を，定数表の値を用いて有効数字3桁で求めよ。
  \end{shoumon}
\notekiroku{2026-10-25 B ◎△◎}
\begin{kotaeutsushi}{2026-10-25}
\toi{1}{◎}{}
\begin{align*}
S&=\frac{3}{2}N\kB-N\kB(\log\lambda^{3}\rho-1)\\
&=N\kB\Bigl(\frac{5}{2}-\log\lambda^{3}\rho\Bigr)
\end{align*}
\[ \frac{S}{N\kB}=\frac{5}{2}-\log\lambda^{3}\rho \]
\toi{2}{△}{計算}
\[ S<0\quad\therefore\ \frac{5}{2}-\log\lambda^{3}\rho<0 \]
\[ \frac{5}{2}<\log\lambda^{3}\rho,\qquad \ee^{5/2}<\lambda^{3}\rho \]
\[ \Bigl(\frac{\ee^{5/2}}{\rho}\Bigr)^{1/3}<\sqrt{\frac{2\pi\hbar^{2}\beta}{m}} \]
（脇で $E=3N/(2\beta)$ $\therefore$ $\beta=3N/(2E)$ を求める）
\[ \Bigl(\frac{\ee^{5/2}}{\rho}\Bigr)^{2/3}<\frac{3\pi N\hbar^{2}}{mE} \]
\[ E<\frac{3\pi N\hbar^{2}}{m}\,\ayamari{\rho^{3/2}\ee^{-3/5}}{\rho^{2/3}\ee^{-5/3}} \]
\naze{前の行の左辺は $\ee^{5/3}\rho^{-2/3}$ なので，$E$ について解くと逆数の $\rho^{2/3}\ee^{-5/3}$ が掛かる。指数の分母と分子を入れ換えてはいけない。正しくは $E<N\ee^{-5/3}(3\pi\hbar^{2}/m)\rho^{2/3}$ である。}
\toi{3}{◎}{}
\[ \lambda^{3}\rho=1\quad\therefore\ \frac{2\pi\hbar^{2}}{\kB T_1m}=\Bigl(\frac{1}{\rho}\Bigr)^{2/3} \]
\begin{align*}
\therefore\ T_1&=\frac{2\pi\hbar^{2}}{\kB m}\rho^{2/3}\\
&=\qty{5.89}{\kelvin}
\end{align*}
\[ \frac{5}{2}=\log\lambda^{3}\rho\quad\therefore\ \lambda^{3}\rho=\ee^{5/2} \]
\[ \frac{2\pi\hbar^{2}}{\kB T_0m}=\Bigl(\frac{\ee^{5/2}}{\rho}\Bigr)^{2/3} \]
\begin{align*}
T_0&=\frac{2\pi\hbar^{2}}{\kB m}\rho^{2/3}\ee^{-5/3}\\
&=\qty{1.11}{\kelvin}
\end{align*}
\end{kotaeutsushi}
\begin{kizuki}
\item $(a/b)^{2/3}$ の逆数は $(b/a)^{2/3}=b^{2/3}a^{-2/3}$ で，指数は $2/3$ のまま符号だけが変わる。$\bigl(\ee^{5/2}/\rho\bigr)^{2/3}=\ee^{5/3}\rho^{-2/3}$ と先に指数を計算してから $E$ について解くと，$\ee^{-5/3}\rho^{2/3}$ が出る。
\item 確かめ方：右辺の次元を見る。$\hbar^{2}/m$ は［エネルギー×長さ$^{2}$］なので，エネルギーにするには［長さ$^{-2}$］の $\rho^{2/3}$ が要る（$\rho^{3/2}$ では長さ$^{-9/2}$ になる）。また (3) の $T_0$ では $\rho^{2/3}\ee^{-5/3}$ と書けていたので，同じ条件の式どうしで形をそろえて見比べる。
\end{kizuki}
\begin{tadashii}
\houshin $S=\kB(\beta E+\log Z_N)$ に与えられた式を入れる。$S<0$ の条件は対数の単調性で $\lambda^{3}\rho$ の条件にし，$\beta=3N/(2E)$ で $E$ の条件に直す。(3) は $\lambda^{3}\rho$ を決まった値にする温度を $\lambda\propto T^{-1/2}$ から解く。
\keisan \shoumonkai{1} $F=-\kB T\log Z_N$ より
\[ S=\frac{E-F}{T}=\kB\bigl(\beta E+\log Z_N\bigr)=\kB\Bigl\{\frac{3N}{2}+N\bigl[1-\log(\lambda^{3}\rho)\bigr]\Bigr\} \]
したがって
\[ \frac{S}{N\kB}=\frac{5}{2}-\log(\lambda^{3}\rho) \]
\shoumonkai{2} $S<0$ は $\log(\lambda^{3}\rho)>5/2$，すなわち $\lambda^{3}\rho>\ee^{5/2}$ と同じである。$E=3N/(2\beta)$ より $\beta=3N/(2E)$ を $\lambda^{2}=2\pi\hbar^{2}\beta/m$ に代入すると $\lambda^{2}=3\pi\hbar^{2}N/(mE)$ になる。条件 $\lambda^{3}\rho>\ee^{5/2}$ の両辺を $2/3$ 乗すると $\lambda^{2}>\ee^{5/3}\rho^{-2/3}$ なので
\[ \frac{3\pi\hbar^{2}N}{mE}>\ee^{5/3}\rho^{-2/3}\quad\Longleftrightarrow\quad E<N\ee^{-5/3}\,\frac{3\pi\hbar^{2}}{m}\,\rho^{2/3} \]
\shoumonkai{3} $\lambda^{3}\rho=1$ は $\lambda^{2}=\rho^{-2/3}$，すなわち $2\pi\hbar^{2}/(m_{\mathrm{He}}\kB T_1)=\rho^{-2/3}$ なので
\[ T_1=\frac{2\pi\hbar^{2}\rho^{2/3}}{m_{\mathrm{He}}\kB}=\frac{\num{6.927e-68}\times\num{7.804e18}}{\num{6.65e-27}\times\num{1.38e-23}}\,\unit{\kelvin}=\frac{\num{5.406e-49}}{\num{9.177e-50}}\,\unit{\kelvin}=\qty{5.89}{\kelvin} \]
ここで $\rho^{2/3}=\qty{7.804e18}{\per\metre\squared}$（$\rho^{1/3}=\qty{2.7935e9}{\per\metre}$ の2乗）を使った。$S=0$ は $\lambda^{3}\rho=\ee^{5/2}$，すなわち $\lambda^{2}=\ee^{5/3}\rho^{-2/3}$ である。$\lambda^{2}\propto1/T$ なので
\[ T_0=\ee^{-5/3}\,T_1=\num{0.1889}\times\qty{5.891}{\kelvin}=\qty{1.11}{\kelvin} \]
\kotae{\shoumonkai{1} $S/(N\kB)=5/2-\log(\lambda^{3}\rho)$ \shoumonkai{2} $\lambda^{3}\rho>\ee^{5/2}$，$E<N\ee^{-5/3}(3\pi\hbar^{2}/m)\rho^{2/3}$ \shoumonkai{3} $T_1=\qty{5.89}{\kelvin}$，$T_0=\qty{1.11}{\kelvin}$}
\kenzan (1) に $\lambda^{3}\rho$ を代入して指数をとると $S=\kB\log[\ee^{5/2}/(\lambda^{3}\rho)]^{N}$ となり，要点整理のエントロピーの式の $[\;]$ の中が $\ee^{5/2}/(\lambda^{3}\rho)$ の $N$ 乗になることと一致する。(2) の $E$ の条件は，要点整理のエントロピーの式で $S<0$ とした条件と一致する。(3) 低温ほど $\lambda$ が大きく $\lambda^{3}\rho$ が大きいので，$T<T_0$ で $S<0$ になり，$T_0<T_1$ である。 実際のヘリウム4は約 \qty{2.17}{\kelvin} で超流動になり，古典的な扱いが破れる温度のオーダーは $T_1$ と合っている。
\tsumazuki $S<0$ になる温度 $T_0$ より高い $T_0<T<T_1$ でも $\lambda^{3}\rho>1$ であり，$S>0$ だからといってこの式が正しいわけではない。$S<0$ は数え間違いが目に見える形で現れた結果である。
\end{tadashii}
\end{mondai}

% toubun: kiroku\toubun\統計力学2\005_2026-10-26.tex
\mondaibangou{10}
\begin{mondai}{005-senyuu}{計算}{1粒子固有状態の占有確率と平均の占有数，$\lambda^{3}/V\ll1$ と $\lambda^{3}\rho\ll1$ の違い，平均の粒子間距離と $\lambda$ の比較}
\begin{teisuu}
  換算プランク定数 & $\hbar$ & \qty{1.05e-34}{\joule\second}\\
  ボルツマン定数 & $\kB$ & \qty{1.38e-23}{\joule\per\kelvin}\\
  ヘリウム4原子の質量 & $m_{\mathrm{He}}$ &\qty{6.65e-27}{\kilogram}\\
  電子の質量 & $\melec$ & \qty{9.11e-31}{\kilogram}\\
\end{teisuu}
体積 $V$ の立方体の容器の中の $N$ 個の同種粒子（相互作用なし）を $Z_N=Z_1^{N}/N!$ で扱う。$\beta=1/(\kB T)$，$\lambda=\sqrt{2\pi\hbar^{2}\beta/m}$，$\rho=N/V$ とする。
  \begin{shoumon}
    \item この扱いでは，1個の粒子がある1粒子固有状態 $\alpha=(n_x,n_y,n_z)$ をとる確率は $p_\alpha=\ee^{-\beta E_\alpha}/Z_1$ である。$\beta E_0\ll1$ のとき $p_\alpha<\lambda^{3}/V$ を示し，状態 $\alpha$ の平均の占有数 $Np_\alpha$ が $\lambda^{3}\rho$ より小さいことを示せ。
    \item 銅の伝導電子を，相互作用もスピンも無い同種粒子の気体とみなす。数密度 $\rho=\qty{8.50e28}{\per\metre\cubed}$，$T=\qty{300}{\kelvin}$ として，$\lambda$ と $\lambda^{3}\rho$ を有効数字3桁で求め，古典的に扱ってよいかを判定せよ。また，試料を一辺 $\qty{1.00}{\centi\metre}$ の立方体として $\lambda^{3}/V$ を求め，和を積分に置き換える仮定 $\beta E_0\ll1$ が成り立つかを述べよ。
    \item (2) の電子の気体で $\lambda^{3}\rho=1$ となる温度 $T_1$ を有効数字3桁で求めよ。
    \item 平均の粒子間距離 $d=\rho^{-1/3}$ を用いて $\lambda^{3}\rho$ を表し，$\lambda^{3}\rho\gg1$ が粒子の波動関数の描像で何を意味するかを説明せよ。また (2) の電子について $d$ を有効数字3桁で求め，$\lambda$ と比べよ。
  \end{shoumon}
\notekiroku{2026-10-26 B ○○◎○}
\begin{kotaeutsushi}{2026-10-26}
\toi{1}{○}{記述}
$\beta E_0\ll1$ のとき $Z_1=V/\lambda^{3}$。

$\ayamari{0\le\beta E_\alpha}{E_\alpha\ge3E_0>0}$ より $\ee^{-\beta E_\alpha}<1$
\naze{$0\le\beta E_\alpha$ からは $\ee^{-\beta E_\alpha}\le1$ しか出ず，強い不等号 $<1$ の根拠にならない。量子数は $1$ 以上なので $E_\alpha\ge3E_0>0$ であり，$\beta>0$ から $\ee^{-\beta E_\alpha}<1$ となる。}
\begin{align*}
p_\alpha&=\ee^{-\beta E_\alpha}\frac{\lambda^{3}}{V}<\frac{\lambda^{3}}{V}\\
Np_\alpha&<\lambda^{3}\frac{N}{V}=\lambda^{3}\rho
\end{align*}
\toi{2}{○}{単位}
\[ \lambda=\num{4.29e-9}\,\nuke{\unit{\metre}} \]
\naze{$\lambda$ は長さなので単位 $\unit{\metre}$ が要る。$\lambda=\qty{4.29e-9}{\metre}$ である。}
\[ \lambda^{3}\rho=\num{6.69e3}\gg1 \]
古典的に扱えない。
\[ \frac{\lambda^{3}}{V}=\num{7.87e-20} \]
\begin{align*}
\beta E_0&=\frac{\pi}{4}\Bigl(\frac{\lambda}{L}\Bigr)^{2}\\
&=\frac{\pi}{4}\Bigl(\frac{\lambda^{3}}{V}\Bigr)^{2/3}\\
&=\num{1.44e-13}
\end{align*}
成り立つ。
\toi{3}{◎}{}
\[ \lambda^{3}\rho=1 \]
\[ \frac{2\pi\hbar^{2}}{\kB Tm}=\Bigl(\frac{1}{\rho}\Bigr)^{2/3} \]
\begin{align*}
T&=\frac{2\pi\hbar^{2}}{\kB m}\rho^{2/3}\\
&=\qty{1.07e5}{\kelvin}
\end{align*}
\toi{4}{○}{記述}
\[ \lambda^{3}\rho=\Bigl(\frac{\lambda}{d}\Bigr)^{3} \]
（図：幅 $\lambda$ に広がった波と，その下に間隔 $d$ をおいて並ぶ波）
\[ d=\num{2.27e-10}\,\nuke{\unit{\metre}} \]
\naze{$\lambda^{3}\rho\gg1$ は $\lambda\gg d$ のことで，各粒子の波動関数の広がり $\lambda$ が隣の粒子までの距離より大きく，波動関数どうしが重なって粒子の統計性が現れることを意味する。図だけでなくこの説明と，$\lambda$ との比較 $\lambda/d\simeq19$ を書く必要がある。$d$ も長さなので単位 $\unit{\metre}$ が要る。}
\end{kotaeutsushi}
\begin{kizuki}
\item (1) の強い不等号の根拠：$E_\alpha\ge0$ だけでは $\ee^{-\beta E_\alpha}\le1$ しか言えない。量子数が $1$ 以上なので $E_\alpha=E_0(n_x^{2}+n_y^{2}+n_z^{2})\ge3E_0>0$ であり，$\beta>0$ と合わせて $\ee^{-\beta E_\alpha}<1$ となる。等号が入るかどうかを，最小の量子数を入れて確かめる。
\item (4) の「描像で何を意味するか」は，$\lambda^{3}\rho=(\lambda/d)^{3}$ の式だけでは答えにならない。$\lambda^{3}\rho\gg1\iff\lambda\gg d$ で，広がり $\lambda$ の波動関数が隣の粒子まで重なるので統計性が効く，と文で書く。「比べよ」と問われたら比 $\lambda/d\simeq19$ を数値で示す。
\item \mae{第1回 問題2・10-25} (2) の $\lambda$ と (4) の $d$ に単位が無い。前回も $\lambda$ の単位を落としている。数値を書いたらその場で次元を確かめ，$\lambda$・$d$ は長さ（$\unit{\metre}$），$\lambda^{3}\rho$・$\lambda^{3}/V$・$\beta E_0$ は無次元，と区別してから答えを書く。
\end{kizuki}
\begin{tadashii}
\houshin (1) はボルツマン因子が $1$ より小さいことと $Z_1=V/\lambda^{3}$ を使う。(2)(3) はヘリウム4の熱的ドブロイ波長の問題と同様に数値を代入し，$\lambda^{3}\rho\propto T^{-3/2}$ を使う。(4) は $\rho=d^{-3}$ で $\lambda^{3}\rho$ を長さの比で表す。
\keisan \shoumonkai{1} 量子数は $1$ 以上なので $E_\alpha\ge3E_0>0$，したがって $\ee^{-\beta E_\alpha}<1$ である。$\beta E_0\ll1$ では $Z_1=V/\lambda^{3}$ なので
\[ p_\alpha=\frac{\ee^{-\beta E_\alpha}}{Z_1}<\frac{1}{Z_1}=\frac{\lambda^{3}}{V} \]
各粒子は独立に確率 $p_\alpha$ で状態 $\alpha$ をとるので，状態 $\alpha$ の平均の占有数は $Np_\alpha<N\lambda^{3}/V=\lambda^{3}\rho$ である。
\shoumonkai{2} 電子の質量を使うと
\[ \lambda^{2}=\frac{2\pi\hbar^{2}}{\melec\kB T}=\frac{\num{6.927e-68}}{\num{9.11e-31}\times\num{1.38e-23}\times\num{300}}\,\unit{\metre\squared}=\frac{\num{6.927e-68}}{\num{3.772e-51}}\,\unit{\metre\squared}=\qty{1.837e-17}{\metre\squared} \]
なので $\lambda=\qty{4.286e-9}{\metre}$，$\lambda^{3}=\qty{7.872e-26}{\metre\cubed}$ である。よって
\[ \lambda^{3}\rho=\num{7.872e-26}\times\num{8.50e28}=\num{6.69e3}\gg1 \]
であり，古典的に扱うことはできない（量子統計が必要）。一方，$V=\qty{1.00e-6}{\metre\cubed}$ より
\[ \frac{\lambda^{3}}{V}=\frac{\num{7.872e-26}}{\num{1.00e-6}}=\num{7.87e-20}\ll1 \]
なので $\beta E_0\ll1$（和を積分に置き換える仮定）は成り立っている。$\lambda^{3}\rho=N\lambda^{3}/V$ なので，$\lambda^{3}/V\ll1$ でも $N$ が大きければ $\lambda^{3}\rho\gg1$ となる。2つの条件は別物である。
\shoumonkai{3} $\lambda^{3}\propto T^{-3/2}$ なので $\lambda^{3}\rho$ は $T^{-3/2}$ に比例する。$T=\qty{300}{\kelvin}$ で $\lambda^{3}\rho=\num{6691}$ なので，これが $1$ になる温度は
\[ T_1=\qty{300}{\kelvin}\times(\num{6691})^{2/3}=\qty{300}{\kelvin}\times\num{355.1}=\qty{1.07e5}{\kelvin} \]
\shoumonkai{4} $\rho=d^{-3}$ より $\lambda^{3}\rho=(\lambda/d)^{3}$ である。$\lambda^{3}\rho\gg1$ は $\lambda\gg d$，すなわち各粒子の波動関数の広がり $\lambda$ が隣の粒子までの距離より大きいことを意味する。このとき粒子の波動関数どうしが重なり，粒子の統計性が現れる。電子については
\[ d=\rho^{-1/3}=(\num{8.50e28})^{-1/3}\,\unit{\metre}=\qty{2.27e-10}{\metre} \]
で，$\lambda/d\simeq\num{18.8}$ である。$\lambda$ は $d$ の約19倍で，波動関数は強く重なっている（丸めずに計算した $\lambda/d=\num{18.84}$ の3乗 $\num{6.69e3}$ は (2) の $\lambda^{3}\rho$ と一致する）。
\kotae{\shoumonkai{1} $p_\alpha<1/Z_1=\lambda^{3}/V$，$Np_\alpha<\lambda^{3}\rho$ \shoumonkai{2} $\lambda=\qty{4.29e-9}{\metre}$，$\lambda^{3}\rho=\num{6.69e3}\gg1$ で古典的には扱えない。$\lambda^{3}/V=\num{7.87e-20}\ll1$ で $\beta E_0\ll1$ は成り立つ \shoumonkai{3} $T_1=\qty{1.07e5}{\kelvin}$ \shoumonkai{4} $\lambda^{3}\rho=(\lambda/d)^{3}$，$d=\qty{2.27e-10}{\metre}$，$\lambda\simeq19d$ で波動関数が重なる}
\begin{ronten}
  \item \shoumonkai{1} 量子数が $1$ 以上なので $E_\alpha\ge3E_0>0$ であり，そこから $\ee^{-\beta E_\alpha}<1$ が言えることを根拠として書く（$E_\alpha\ge0$ だけでは不等号が厳密にならない）。
  \item \shoumonkai{1} $Z_1=V/\lambda^{3}$ を使うには $\beta E_0\ll1$ が必要なことと，各粒子が独立に確率 $p_\alpha$ で状態 $\alpha$ をとるので平均の占有数が $Np_\alpha$ になることを述べる。
  \item \shoumonkai{2} $\lambda^{3}/V\ll1$（和を積分にする条件）と $\lambda^{3}\rho\ll1$（古典的な数え方の条件）が別の条件であることを区別して判定する。
  \item \shoumonkai{4} $\lambda^{3}\rho=(\lambda/d)^{3}$ なので，$\lambda^{3}\rho\gg1$ は広がり $\lambda$ の波動関数が隣の粒子まで重なることを意味すると説明する。
\end{ronten}
\kenzan $\lambda\propto m^{-1/2}$ なので，ヘリウム4（熱的ドブロイ波長の問題）との比は $\lambda_{\mathrm{e}}/\lambda_{\mathrm{He}}=\sqrt{m_{\mathrm{He}}/\melec}=\sqrt{\num{7300}}\simeq\num{85.4}$ で，$\num{5.016e-11}\times\num{85.4}\simeq\num{4.29e-9}$ と一致する。$T_1$ を $\lambda^{3}\rho=1$ の定義式 $T_1=2\pi\hbar^{2}\rho^{2/3}/(\melec\kB)$ で直接計算しても $\qty{1.07e5}{\kelvin}$ になる。 実際の伝導電子はスピン $1/2$ をもちフェルミ統計に従う（第9章で扱う）が，室温で古典的に扱えないという結論は変わらない。
\tsumazuki $\lambda^{3}/V\ll1$（積分近似の条件）と $\lambda^{3}\rho\ll1$（古典的な数え方の条件）を混同しないこと。前者が成り立っていても後者が破れることがある。
\end{tadashii}
\end{mondai}

\noteowari
\end{document}
